\section{緒言}

文字入力は，情報端末を通じて人と人，人と社会をつなぐ最も基本的な行為の一つである．
しかし加齢に伴う運動・視機能の変化，身体的制約，端末や状況によって，
同じ入力操作が大きな負担となることがある．
誰もが負担少なく自分の意思を伝えられることは，生活の質と尊厳を支えるうえで重要であり，
本研究はこうした人間の負担に着目した使いやすい入力支援を目指す．

従来の日本語入力方式（IME）は十分に実用的かつ高性能である．
本研究の問題意識は，IME の変換精度ではなく，
ユーザに要求される\textbf{入力操作量}にある．
すなわち課題は「変換が正しくない」ことではなく，
「人間の側に求められる操作が多い」ことにある．
かな入力は 50 音相当，QWERTY は 26 キー以上，フリック入力は 12 キーと方向操作・習熟を要求する．
これに伴い，小さなキーを正確に押し分ける負担，フリック方向の記憶・習熟の負担，
多数の候補から目的語を探す認知的負荷，多くのキーが画面を占有する視認性低下が生じ，
小型端末，高齢者，身体制約者などでとりわけ重い負担となりうる．
実際，高齢者ではタッチ操作の「文字数字入力」などで加齢に伴う成功率の低下が報告されており\cite{hatakenaka2024touch}，
認知機能の個人差が Web 操作の難しさに影響することも指摘されている\cite{suzuki2009elderly}．

本研究では，入力情報を日本語の 5 母音（\texttt{a, i, u, e, o}）および撥音
（\texttt{n}）にまで削減する．
たとえば「ありがとう」は \texttt{a i a o u} と入力される．
入力に用いるキーを 6 つにまで絞ることで，
スマートフォンの画面においてキーボードが占める割合を大幅に減らすことができ，
一つひとつのキーを大きく配置できる．
これは，小さなキーを正確に押し分けることが負担となりやすい高齢者にとって特に有効と考えられる．
一方で，このとき子音情報や語彙情報の大部分が失われるため，
入力は本質的に曖昧（ファジィ）となり，母音列は一つの文に一対一で対応せず複数の候補へ多義的に広がる． % 追加: ファジィ入力の位置づけを一文に集約
そこで失われた情報を LLM による CTX 推定で補い，
\textbf{母音列からの日本語候補生成・順位付け}として扱う．
これは，正確な読みを人間が打ち込み候補を探す方式から，曖昧な母音IMEを LLM がリランキングして提示する方式への発想の転換であり，
元の文章の完全復元ではなく候補提示型の入力支援として位置づける．
